%!TEX root=../main.tex

\section{Equivalence of ReLU and Gated ReLU}\label{sec:relu-grelu-equivalence}

This section builds upon our convex re-formulations to show that the full C-ReLU problem is equivalent to a sub-sampled C-GReLU problem up to the norm of their optimal solutions.
As a consequence, we give an approximation algorithm for ReLU networks that first computes the solution to C-GReLU and then solves an auxiliary \emph{cone decomposition} problem.
The cone decomposition can be formulated as a linear program (LP) or second-order cone program (SOCP), and admits a closed form solution when \( X \) is full row-rank.
Before presenting these fully-general results, we study the unregularized setting, where we show the approximation is exact.
All proofs are deferred to \cref{app:relu-grelu-equivalence}.

\begin{figure}
    \centering
    \input{figures/chapter_two/cone_decomp}
    \caption{An illustration of the Cone Decomposition (CD) procedure: \( u_i \) is decomposed onto the Minkowski difference \( \calK_i - \calK_i \).}%
    \label{fig:cone-decomp}
    \vspace{-0.3cm}
\end{figure}

Let \( \lambda = 0 \) and consider the C-GReLU problem~\eqref{convex-forms:eq:unconstrained-relu-mlp}.
For each \( D_i \in \tilde \calD \), we seek to decompose the optimal data-local models as \( u^*_i = v_i - w_i \in \calK_i - \calK_i \).
If these decompositions exist, collecting them into \( (v, w) = \cbr{(v_i, w_i)} \) gives a feasible point for the C-ReLU problem with the same optimal objective value as C-GReLU.
The next proposition gives sufficient conditions on the data for this to happen.
\begin{restatable}{proposition}{nonSingularCones}\label{prop:non-singular-cones}
    If \( X \) is full row-rank, then \( \calK_i - \calK_i = \R^d  \) for every \( D_i \in \calD_{\calX} \).
    As a result, the C-ReLU, C-GReLU, NC-ReLU, and NC-GReLU problems are all equivalent.
\end{restatable}
Unfortunately, \cref{prop:non-singular-cones} does not extended to \( n > d \);
in \cref{prop:col-rank-counter-example}, we give full-rank \( X \) for which some \( \calK_i \) is contained in a subspace of \( \R^d \), implying \( \calK_i - \calK_i \subset \R^d \).
We call such cones (and associated gate vectors) \emph{singular}.
\begin{restatable}{proposition}{coneContainment}\label{prop:cone-containment}
    Suppose \( \calK_i \) is singular for \( D_i \in \calD_{X} \).
    Then \( \exists D_j \in \calD_{X} \) such that \( \calK_j - \calK_j = \R^d \) and \( \calK_i \subset \calK_j \).
\end{restatable}
That is, every singular cone is contained within a non-singular cone.
As a result, we show that these ``bad'' cones,
\begin{equation}\label{eq:singular-cones}
    \calS(\tilde \calD_X) = \cbr{ D_i \in \tilde \calD_X : \calK_i - \calK_i \subset \R^d },
\end{equation}
can be safely ignored when forming the convex programs.
\begin{restatable}{theorem}{coneElimination}\label{thm:cone-elimination}
    Let \( \tilde \calD \subseteq \calD_X \) and \( \lambda \geq 0 \).
    Then the C-ReLU problem with \( \tilde \calD \) is equivalent
    to the C-ReLU problem with \( \tilde \calD \setminus \calS(\tilde \calD) \).
    If \( \lambda = 0 \), then both problems are equivalent to the sub-sampled
    C-GReLU problem with \( \tilde \calD \setminus \calS(\tilde \calD) \).
    %That is, the same sub-sampled problems without singular cones.
\end{restatable}
%In words, \cref{alg:cone-decomp} says that singular cones can be safely omitted
%and that the gated ReLU and ReLU problems are equivalent we need only omit all singular cones
%for the C-ReLU and C-GReLU problems to be equivalent.
Choosing \( \tilde \calD = \calD_X \) shows that the full C-ReLU
problem is exactly equivalent to the unconstrained C-GReLU problem without
singular gate vectors.
\cref{alg:cone-decomp} provides a template for training ReLU networks by leveraging cone decompositions and \cref{thm:cone-elimination}.
Note that \( \tilde \calD \) is generated by randomly sampling gate vectors \( g_i \sim \calN(0, I) \).
This is sufficient to recover the C-ReLU problem as \cref{thm:cone-elimination} implies singular cones, for which the sampling probability is zero, don't contribute to the solution.

%% MP:This figure looks great!
\subsection{Approximating ReLU by Cone Decompositions}

We have seen that decomposing \( u_i^* = v_i - w_i \) allows us to map the C-GReLU problem into the C-ReLU problem.
However, triangle inequality shows \( \norm{u_i^*}_2 \leq \norm{v_i}_2 + \norm{w_i}_2 \), meaning the cone decomposition can only increase the norm of the model (see \cref{fig:cone-decomp}).
For \( \lambda > 0 \), this increases the penalty term in the objective (Eq.~\ref{convex-forms:eq:convex-relu-mlp}), although the loss \( \calL \) is unchanged.
This section develops cone decomposition algorithms for which we know the blow-up of the norm is not too large.
As a result, we obtain approximation guarantees for solving C-ReLU by solving C-GReLU.

In what follows, \( \calK = \cbr{w : (2 D - I) X w \succeq 0} \) denotes a non-singular cone, \( \tilde X = (2 D - I) X \), and \( \kappa(A) \) is the ratio of the largest and smallest \emph{non-zero} singular values of \( A \).
Our first result gives conditions for the existence of a closed-form
decomposition.
\begin{restatable}{proposition}{closedFormDecomp}\label{prop:closed-form-decomp}
    Suppose \( X \) is full row-rank.
    If \( \calI = \cbr{i \in [n] : \abr{\tilde x_i, u} < 0 } \), then for every \( u \in \R^d \),
    \vspace{-1ex}
    \[
        u = \rbr{u + w} - w, \text{where} \,  w = - \tilde X_{\calI}^\dagger \tilde X_{\calI} u,
    \]
    \vspace{-1ex}
    is a valid decomposition onto \( \calK - \calK \) satisfying,
    \[
        \norm{u + w}_2 + \norm{w}_2 \leq 2 \norm{u}_2.
    \]
\end{restatable}
In general, we cannot hope for constant approximations since \( n \gg d \) implies
the cones \( \calK \) are very ``narrow''.
\begin{restatable}{proposition}{worstCaseNorm}\label{prop:worst-case-norm}
    There does not exist a decomposition \( u = v - w \), where \( v, w \in \calK \), such that
    \[
        \norm{v}_2 + \norm{w}_2 \leq C \norm{u},
    \]
    holds for an absolute constant \( C \).
\end{restatable}
\begin{algorithm}[tb]
    \caption{Solving C-ReLU by Cone Decomposition}
    \label{alg:cone-decomp}
    \begin{algorithmic}
        \STATE {\bfseries Input:} data \( (X, y) \), \( \lambda \geq 0 \), num. samples \( p \), objective \( R \).
        \STATE \textbf{Sample}: \( \tilde \calD \!=\! \cbr{\diag(\mathds{1}(X g_i \geq 0)) :  g_i \!\sim \!\calN(0, I), i \in [p] } \)
        \STATE \textbf{Solve C-GReLU}:
        \STATE \hspace{1em} \( u^* \in \argmin_{u} \calL(\sum_{\tilde \calD} D_i X u_i, y) + \lambda \sum_{\tilde \calD} \norm{u_i}_2\)
        \STATE \textbf{Solve Cone Decomposition}:
        \STATE \hspace{1em} \( \bar v, \bar w \in \argmin_{v, w} \cbr{ R(v, w) : u_i^* = v_i - w_i, i \in [p] } \)
        \STATE \textbf{Return}: \((\bar v, \bar w)\)
    \end{algorithmic}
\end{algorithm}
\vspace{-1ex}
When \( X \) is not full row-rank, we can solve
\begin{equation}\label{eq:decomposition-program}
    \begin{aligned}
        \textbf{CD}: \quad \min_{v, w \in \calK} \cbr{ R(v, w) : v - w = u },
    \end{aligned}
\end{equation}
where \( R : \R^{d \times d} \mapsto \R \) is some loss function.
Taking \( R(v, w) = 0 \) reduces to a linear feasibility problem which can be handled by off-the-shelf LP solvers.
Choosing \( R(v,w) = \norm{v}_2 + \norm{w}_2 \) yields a second-order cone program (SOCP) for which we have the following guarantee.
\begin{restatable}{proposition}{socpDecomp}\label{prop:socp-decomp}
    For every \( u \in \R^d \), if \( \rbr{\bar v, \bar w} \) is a solution to the cone-decomposition program~\eqref{eq:decomposition-program} with
    \( R(v,w) = \norm{v}_2 + \norm{w}_2 \), then there exists \( \calJ \subseteq [n] \) such that
    \[
        \norm{\bar v}_2 + \norm{\bar w}_2 \leq \rbr{1 + 2 \kappa(\tilde X_{\calJ})} \norm{u}_2.
    \]
\end{restatable}

Note that the general setting incurs a penalty of \( \kappa(\tilde X_{\calJ}) \) compared to \cref{prop:closed-form-decomp}.
Intuitively, this term measures the narrowness of \( \calK \) and the difficulty of the decomposition.
Combining \cref{prop:socp-decomp} with \cref{thm:cone-elimination} gives our main approximation result.
\begin{restatable}{theorem}{approxResult}\label{thm:approx-result}
    Let \( \lambda \geq 0 \) and let \( p^* \) be the optimal value of the full
    C-ReLU problem with training set \( (X, y) \).
    There exists \( \calJ \subseteq [n] \) such that the C-GReLU problem with
    patterns \( \calD_X \setminus \calS(\calD_X) \), minimizer \( u^* \),
    and optimal value \( d^* \) satisfies,
    \[
        d^*
        \leq p^*
        \leq d^*
        + 2 \lambda \kappa(\tilde X_{\calJ}) \sum_{D_i \in \tilde \calD} \norm{u_i^*}_2.
    \]
\end{restatable}
As a consequence of \cref{thm:approx-result}, \cref{alg:cone-decomp} is guaranteed to approximate the C-ReLU problem if \( R(v,w) = \norm{v}_2 + \norm{w}_2 \) and \( p \) is sufficiently large.
As \( \lambda \into 0 \), this result smoothly recovers \cref{thm:cone-elimination}, implying we can control the approximation by adjusting the regularization.

\begin{figure}
    \centering
    \input{figures/chapter_two/relations}
    \caption{Summary of equivalences between convex (blue) and non-convex (red) neural
        network training problems with gated ReLU (left) and ReLU (right)
        activations.
        The convex programs C-GReLU and C-ReLU are equivalent to the standard
        non-convex training problems NC-GReLU and NC-GReLU and are related to
        each other via the cone decomposition procedure.
    }%
    \label{fig:relations}
    \vspace{-0.3cm}
\end{figure}
